\section{总结与延伸}

\subsection{讲者的核心总结}

Justin Johnson 在课程结尾强调了两个核心贡献：

\begin{enumerate}
    \item \textbf{Attention}：一种新的计算原语，本质上操作于向量集合之上。它高度可并行、可扩展，核心计算只有几次矩阵乘法。
    \item \textbf{Transformer}：以 Self Attention 为核心的神经网络架构。自 2017 年以来统治了深度学习的几乎所有领域，从语言到视觉到多模态。
\end{enumerate}

\subsection{全课知识图谱}

\begin{center}
\begin{tikzpicture}[node distance=0.8cm, every node/.style={draw, rounded corners, minimum width=3.5cm, minimum height=0.7cm, font=\small}]
\node (rnn) {RNN Encoder-Decoder};
\node (bottle) [below=of rnn] {固定向量瓶颈};
\node (att) [below=of bottle] {Attention 机制};
\node (gen) [below=of att] {通用注意力算子};
\node (self) [below=of gen] {Self Attention};
\node (tf) [below=of self] {Transformer};
\draw[->, thick] (rnn) -- (bottle);
\draw[->, thick] (bottle) -- (att);
\draw[->, thick] (att) -- (gen);
\draw[->, thick] (gen) -- (self);
\draw[->, thick] (self) -- (tf);
\node[draw=none, right=0.8cm of rnn, text width=5cm, font=\scriptsize] {Seq2Seq 机器翻译};
\node[draw=none, right=0.8cm of bottle, text width=5cm, font=\scriptsize] {Context Vector 信息丢失};
\node[draw=none, right=0.8cm of att, text width=5cm, font=\scriptsize] {动态回看输入、可微对齐};
\node[draw=none, right=0.8cm of gen, text width=5cm, font=\scriptsize] {缩放点积、KQV 分离、矩阵乘法};
\node[draw=none, right=0.8cm of self, text width=5cm, font=\scriptsize] {置换等变、位置编码、Masking};
\node[draw=none, right=0.8cm of tf, text width=5cm, font=\scriptsize] {MHA + MLP + 残差 + LayerNorm};
\end{tikzpicture}
\end{center}

\subsection{关键 Takeaways}

\begin{importantbox}{六条核心原则}
\begin{enumerate}
    \item \textbf{注意力的本质是动态的信息检索}：给定查询，从数据中按相关性加权提取信息
    \item \textbf{KQV 分离是关键设计}：让网络分别学习``如何匹配''（Key-Query）和``返回什么''（Value）
    \item \textbf{Self Attention 操作于集合}：天然不感知顺序，需要位置编码来注入序列信息
    \item \textbf{注意力本质上是矩阵乘法}：四次 matmul 就是整个 Multi-Head Self Attention
    \item \textbf{Transformer = Attention + MLP + 残差 + 归一化}：简洁的组合，强大的效果
    \item \textbf{并行性决定可扩展性}：Transformer 之所以胜出，是因为它能充分利用大规模并行硬件
\end{enumerate}
\end{importantbox}

\end{document}
